% AA-23 | Algebra abstrakcyjna | Iteracja XII
\chapter{Ślad, norma i wyróżnik w rozszerzeniach ciał}
\label{aa:ch:trace-norm-discriminant}
\label{aa:ch:trace-norm}

\section{Od mnożenia w ciele do operatora liniowego}
W Rozdziale~\ref{aa:ch:galois-correspondence} opisaliśmy grupy automorfizmów
rozszerzeń Galois. Innym sposobem wydobywania informacji z rozszerzenia
jest potraktowanie mnożenia przez ustalony element jako przekształcenia
liniowego. Otrzymamy w ten sposób \emph{ślad} i \emph{normę}.
Konstrukcja ma sens również wtedy, gdy rozszerzenie nie jest rozdzielcze.

\begin{definicja}[Ślad i norma elementu rozszerzenia]
\label{aa:def:field-trace-norm}
Niech $L/K$ będzie rozszerzeniem ciał skończonego stopnia $n=[L:K]$.
Dla $a\in L$ niech $m_a:L\to L$ oznacza $K$-liniowy endomorfizm
$m_a(x)=ax$. Definiujemy
\[
\operatorname{Tr}_{L/K}(a)=\operatorname{tr}_K(m_a),\qquad
N_{L/K}(a)=\det_K(m_a).
\]
Wyznacznik i ślad po prawej stronie obliczamy w dowolnej bazie
$K$-przestrzeni $L$. Oba należą do $K$.
Ponadto wielomianem charakterystycznym elementu $a$
względem $L/K$ nazywamy $\chi_{a,L/K}(T)=\det(TI-m_a)\in K[T]$.
\end{definicja}

\begin{uwaga}[Zależność od rozszerzenia]\label{aa:uw:trace-relative-extension}
Indeksy $L/K$ są częścią definicji. Jeżeli $a\in F$ dla ciała pośredniego
$K\subseteq F\subseteq L$, to operatory mnożenia przez $a$
na $F$ i na $L$ działają na przestrzeniach o różnych wymiarach;
wobec tego ich ślady i wyznaczniki nie muszą być równe.
Wielomian $\chi_{a,L/K}$ ma zawsze stopień $[L:K]$,
niezależnie od stopnia minimalnego wielomianu $a$ nad $K$.
Przypadek $L=K$ jest również objęty definicją:
wtedy $m_a$ ma macierz $(a)$, a więc
$\operatorname{Tr}_{K/K}(a)=N_{K/K}(a)=a$.
\end{uwaga}

\begin{lemat}[Niezależność od wyboru bazy]
\label{aa:lem:trace-norm-basis}
Niech $L/K$ będzie skończonym rozszerzeniem ciał, a $a\in L$.
Ślad $\operatorname{Tr}_{L/K}(a)$, norma $N_{L/K}(a)$
oraz wielomian charakterystyczny $\chi_{a,L/K}$
nie zależą od wyboru bazy $K$-przestrzeni $L$.
\end{lemat}
\begin{proof}
Jeżeli $A$ i $B$ są macierzami tego samego endomorfizmu $m_a$
w dwóch bazach, to istnieje odwracalna macierz $S$ taka, że
$B=S^{-1}AS$. Stąd $\operatorname{tr}B=\operatorname{tr}A$,
gdyż $\operatorname{tr}(UV)=\sum_{i,j}u_{ij}v_{ji}=\operatorname{tr}(VU)$,
oraz $\det B=(\det S)^{-1}(\det A)\det S=\det A$.
Ponadto $TI-B=S^{-1}(TI-A)S$, a więc
$\det(TI-B)=\det(TI-A)$ w $K[T]$.
Zatem ślad, norma i wielomian charakterystyczny
nie zależą od wyboru bazy.
\end{proof}

\begin{twierdzenie}[Podstawowe własności śladu i normy]
\label{aa:tw:trace-norm-basic}
\label{aa:thm:trace-norm-properties}
Niech $L/K$ będzie skończonym rozszerzeniem stopnia $n$.
Dla wszystkich $a,b\in L$ i $c\in K$ zachodzą
\begin{enumerate}[label=\textup{(\roman*)}]
\item $\operatorname{Tr}_{L/K}(a+b)=\operatorname{Tr}_{L/K}(a)+\operatorname{Tr}_{L/K}(b)$;
\item $\operatorname{Tr}_{L/K}(ca)=c\operatorname{Tr}_{L/K}(a)$;
\item $N_{L/K}(ab)=N_{L/K}(a)N_{L/K}(b)$;
\item $\operatorname{Tr}_{L/K}(c)=nc$ oraz $N_{L/K}(c)=c^n$;
\item $N_{L/K}(a)\ne0$ wtedy i tylko wtedy, gdy $a\ne0$;
      w szczególności $N_{L/K}(a^{-1})=N_{L/K}(a)^{-1}$ dla $a\ne0$.
\end{enumerate}
\end{twierdzenie}
\begin{proof}
Mamy $m_{a+b}=m_a+m_b$ oraz $m_{ca}=c\,m_a$.
Liniowość śladu daje (i) i (ii). Z łączności mnożenia
w $L$ wynika $m_{ab}=m_a\circ m_b$, więc multiplikatywność
wyznacznika daje (iii). Dla $c\in K$ endomorfizm $m_c$
ma w każdej bazie macierz $cI_n$, skąd (iv).
Jeśli $a\ne0$, to mnożenie przez $a$ ma odwrotność $m_{a^{-1}}$,
zatem jego wyznacznik jest niezerowy. Jeśli $a=0$, to $m_0=0$,
więc jego wyznacznik jest zerem. Ostatnia równość wynika z (iii)
i $N_{L/K}(1)=1$.
\end{proof}

\begin{uwaga}[Ślad nie jest homomorfizmem pierścieni]
\label{aa:uw:trace-not-ring}
Ślad jest $K$-liniowy, lecz zazwyczaj nie jest multiplikatywny.
Norma jest multiplikatywna, lecz zazwyczaj nie jest addytywna.
Na przykład dla $L=\mathbb Q(\sqrt2)$ mamy
$\operatorname{Tr}_{L/\mathbb Q}(1)=2$ i
$\operatorname{Tr}_{L/\mathbb Q}(1)^2=4$,
podczas gdy $\operatorname{Tr}_{L/\mathbb Q}(1\cdot1)=2$.
Ponadto $N_{L/\mathbb Q}(1+1)=4$, ale
$N_{L/\mathbb Q}(1)+N_{L/\mathbb Q}(1)=2$.
\end{uwaga}

\section{Macierz mnożenia i wielomian minimalny}
\begin{twierdzenie}[Ślad i norma pierwiastka wielomianu minimalnego]
\label{aa:tw:trace-norm-companion}
Niech $L=K(\alpha)$, gdzie $\alpha$ jest algebraiczny nad $K$,
a jego unormowany wielomian minimalny ma postać
$m_{\alpha,K}(X)=X^n+a_{n-1}X^{n-1}+\cdots+a_1X+a_0$.
Wtedy
\[
\operatorname{Tr}_{L/K}(\alpha)=-a_{n-1},\qquad
N_{L/K}(\alpha)=(-1)^n a_0.
\]
\end{twierdzenie}
\begin{proof}
Na mocy Twierdzenia~\ref{aa:tw:simple-algebraic-extension}
układ $(1,\alpha,\ldots,\alpha^{n-1})$ jest bazą $L$ nad $K$.
Mnożenie przez $\alpha$ przesuwa każdy z pierwszych $n-1$
elementów bazy na następny, a z równości $m_{\alpha,K}(\alpha)=0$
otrzymujemy
$\alpha^n=-a_0-a_1\alpha-\cdots-a_{n-1}\alpha^{n-1}$.
Dla $n=1$ mamy $L=K$, $\alpha=-a_0$ i macierz
$m_\alpha$ jest równa $(-a_0)$; obie formuły
wynikają zatem bezpośrednio z definicji.
Dla $n\ge2$ macierz $m_\alpha$ ma postać blokową
\[
C=\begin{pmatrix}
0_{1\times(n-1)}&-a_0\\
I_{n-1}&-\begin{pmatrix}a_1\\\vdots\\a_{n-1}\end{pmatrix}
\end{pmatrix}.
\]
To jest macierz towarzysząca, poprawnie określona także
dla $n=2$. Jej ślad wynosi $-a_{n-1}$.
Rozwijając wyznacznik względem pierwszego wiersza,
możemy wybrać tylko wpis $-a_0$; jego dopełnienie
jest macierzą jednostkową stopnia $n-1$.
Zatem $\det C=(-a_0)(-1)^{1+n}=(-1)^n a_0$.
\end{proof}

\begin{twierdzenie}[Wielomian charakterystyczny elementu]
\label{aa:thm:trace-norm-minpoly}
Niech $L/K$ będzie skończonym rozszerzeniem ciał,
$a\in L$, $d=[K(a):K]$, $r=[L:K(a)]$, a
$m_{a,K}(T)=T^d+c_{d-1}T^{d-1}+\cdots+c_0$
niech oznacza unormowany wielomian minimalny $a$ nad $K$.
Wtedy wielomian charakterystyczny elementu $a$ względem $L/K$
ma postać $\chi_{a,L/K}(T)=m_{a,K}(T)^r$, a więc
\[
\operatorname{Tr}_{L/K}(a)=-r c_{d-1},\qquad
N_{L/K}(a)=\bigl((-1)^d c_0\bigr)^r.
\]
Nie zakładamy rozdzielczości rozszerzenia.
\end{twierdzenie}
\begin{proof}
Wybierzmy bazę $u_1,\ldots,u_r$ przestrzeni $L$ nad $K(a)$.
Każdy element $L$ można jednoznacznie zapisać jako
kombinację liniową nad $K$ elementów
$a^j u_i$, $0\le j<d$, $1\le i\le r$:
istnienie wynika z rozwinięcia współczynników
w bazie $(1,a,\ldots,a^{d-1})$ ciała $K(a)$ nad $K$,
a niezależność wynika z niezależności $u_i$ nad $K(a)$
i potęg $a^j$ nad $K$. Jest to więc baza $L$ nad $K$.
Mnożenie przez $a$ pozostawia każdą podprzestrzeń
$K$ rozpiętą przez $u_i,au_i,\ldots,a^{d-1}u_i$.
Na tej podprzestrzeni jego macierz jest macierzą
towarzyszącą $m_{a,K}$, gdyż $m_{a,K}(a)=0$.
Wielomian charakterystyczny takiego bloku to $m_{a,K}$.
Aby uzasadnić ten wzór, zauważmy, że w macierzy $TI-C$
na przekątnej stoją wpisy $T$, tuż pod nią wpisy $-1$,
a w ostatniej kolumnie wpisy $c_0,\ldots,c_{d-1}$
(z dodanym $T$ na ostatniej pozycji).
W rozwinięciu wyznacznika względem ostatniej kolumny
składnik pochodzący od $c_{i-1}$ wynosi
$c_{i-1}T^{i-1}$: dopełnienie składa się z $i-1$
czynników $T$ i $d-i$ czynników $-1$, których znaki
znoszą znak dopełnienia algebraicznego.
Wpis $T$ z ostatniego wiersza daje dodatkowo $T^d$.
Dla $d=1$ formuła jest bezpośrednia.
Otrzymujemy zatem $\det(TI-C)=T^d+c_{d-1}T^{d-1}+\cdots+c_0$.
Macierz całego $m_a$ jest blokowo diagonalna
z $r$ takimi blokami. Zatem $\chi_{a,L/K}=m_{a,K}^r$.
Ponieważ dla każdej macierzy $n\times n$ współczynnik
przy $T^{n-1}$ w $\det(TI-A)$ jest równy
$-\operatorname{tr}A$, a wyraz stały to
$(-1)^n\det A$, porównanie współczynników
po obu stronach daje obie formuły.
\end{proof}

\begin{przyklad}[Rozszerzenie kwadratowe]
\label{aa:prz:trace-norm-quadratic}
\label{aa:ex:quadratic-norm}
Niech $d\in\mathbb Q$ będzie liczbą, która nie jest kwadratem
liczby wymiernej, i $L=\mathbb Q(\sqrt d)$.
Dla $a,b\in\mathbb Q$ mnożenie przez $z=a+b\sqrt d$
ma w bazie $(1,\sqrt d)$ macierz
\[
\begin{pmatrix}a&bd\\ b&a\end{pmatrix}.
\]
Otrzymujemy dokładnie
$\operatorname{Tr}_{L/\mathbb Q}(a+b\sqrt d)=2a$,
$N_{L/\mathbb Q}(a+b\sqrt d)=a^2-db^2$.
Zatem element $a+b\sqrt d$ jest niezerowy wtedy i tylko wtedy,
gdy $a^2-db^2\ne0$, a jego odwrotność ma postać
$(a-b\sqrt d)/(a^2-db^2)$. Zauważmy, że obliczenie to
nie korzysta z istnienia automorfizmów rozszerzenia.
\end{przyklad}

\begin{przyklad}[Zmiana ciała, w którym obliczamy normę]
\label{aa:prz:trace-norm-relative-biquadratic}
Niech $K=\mathbb Q$, $F=\mathbb Q(\sqrt2)$ oraz
$L=\mathbb Q(\sqrt2,\sqrt3)$. Ponieważ $[F:K]=2$, sprawdźmy, że $\sqrt3\notin F$.
Gdyby $\sqrt3=u+v\sqrt2$ dla $u,v\in\mathbb Q$, to po
podniesieniu do kwadratu dostalibyśmy
$3=u^2+2v^2+2uv\sqrt2$. Wobec niewymierności
$\sqrt2$ musi być $uv=0$. Dla $v=0$ mielibyśmy
$u^2=3$, a dla $u=0$ mielibyśmy $v^2=3/2$;
żadna z tych liczb nie jest kwadratem liczby wymiernej
(co wynika z rozkładów liczników i mianowników na
czynniki pierwsze). Zatem $\sqrt3\notin F$,
$[L:F]=2$ i $[L:K]=4$.
Dla $a=\sqrt2$ wielomian minimalny nad $K$ to $T^2-2$.
Zatem w rozszerzeniu $F/K$ wzory z
Twierdzenia~\ref{aa:tw:trace-norm-companion} dają
$\operatorname{Tr}_{F/K}(a)=0$ i $N_{F/K}(a)=-2$.
Na $L$ operator mnożenia przez $a$ ma w bazie
$(1,\sqrt2,\sqrt3,\sqrt6)$ dwa identyczne bloki
$\left(\begin{smallmatrix}0&2\\1&0\end{smallmatrix}\right)$:
pierwszy działa na $\operatorname{span}_K(1,\sqrt2)$,
a drugi na $\operatorname{span}_K(\sqrt3,\sqrt6)$.
W konsekwencji $\operatorname{Tr}_{L/K}(a)=0+0=0$,
lecz $N_{L/K}(a)=(-2)^2=4$.
Obliczenie pokazuje, że nie wolno usuwać indeksu $L/K$
ani utożsamiać normy elementu w różnych rozszerzeniach.
\end{przyklad}

\begin{przyklad}[Rozszerzenie trzeciego stopnia]
\label{aa:prz:trace-norm-cubic}
Dla $L=\mathbb Q(\alpha)$, $\alpha=\sqrt[3]2$,
minimalny wielomian $\alpha$ wynosi $X^3-2$:
wielomian ten jest nierozkładalny nad $\mathbb Q$
na mocy kryterium Eisensteina z liczbą pierwszą $2$
(Twierdzenie~\ref{aa:tw:eisenstein}).
Dlatego $\operatorname{Tr}_{L/\mathbb Q}(\alpha)=0$
i $N_{L/\mathbb Q}(\alpha)=2$. Równość dotycząca normy
można również otrzymać z
$N(\alpha)^3=N(\alpha^3)=N(2)=2^3$:
ponieważ $N(\alpha)\in\mathbb Q$, z równości
$N(\alpha)^3=8$ wynika $N(\alpha)=2$.
Bezpośredni wzór z wielomianu minimalnego nie wymaga
rozwiązywania takiego równania.
\end{przyklad}

\begin{przyklad}[Zanik śladu w rozszerzeniu nierozdzielczym]
\label{aa:ex:inseparable-trace}
Niech $p$ będzie liczbą pierwszą,
$K=\mathbb F_p(t)$ i $L=K(s)$, gdzie $s^p=t$.
Wielomian $X^p-t\in\mathbb F_p[t][X]$ jest
nierozkładalny nad $K$ na mocy kryterium Eisensteina
względem elementu pierwszego $t$ w pierścieniu
$\mathbb F_p[t]$ i lematu Gaussa.
Zatem $[L:K]=p$.
Dla dowolnego $a\in L$ mamy $a^p\in K$,
ponieważ $K=\mathbb F_p(t)\subseteq\mathbb F_p(s)$,
$s^p=t$ i stąd $L=\mathbb F_p(s)$;
każdy ułamek wymierny w $s$, podniesiony do potęgi $p$,
jest ułamkiem wymiernym w $s^p=t$ o współczynnikach
z $\mathbb F_p$.
Jeśli $a\in K$, to z podstawowych własności śladu
$\operatorname{Tr}_{L/K}(a)=pa=0$.
Jeśli $a\notin K$, to $[K(a):K]$ jest dzielnikiem
liczby pierwszej $[L:K]=p$, więc jest równy $p$.
Wielomian $X^p-a^p\in K[X]$ ma stopień $p$
i zeruje $a$, zatem jest wielomianem minimalnym $a$.
Nie ma współczynnika przy $X^{p-1}$;
z Twierdzenia~\ref{aa:thm:trace-norm-minpoly}
wynika $\operatorname{Tr}_{L/K}(a)=0$.
Zatem cały funkcjonał śladu $\operatorname{Tr}_{L/K}$
jest zerowy, choć $L/K$ ma dodatni stopień.
Przykład pokazuje, dlaczego w twierdzeniach o
niezdegenerowaniu formy śladu wymagamy rozdzielczości.
\end{przyklad}

\begin{uwaga}[Norma w tym samym nierozdzielczym rozszerzeniu]
\label{aa:uw:pure-inseparable-norm}
W poprzednim przykładzie norma nie znika. Dokładniej,
dla każdego $a\in L$ zachodzi $N_{L/K}(a)=a^p\in K$.
Jeśli $a\in K$, jest to równość $N(a)=a^{[L:K]}=a^p$.
Jeśli $a\notin K$, wielomian minimalny ma postać $X^p-a^p$.
Wzór z Twierdzenia~\ref{aa:thm:trace-norm-minpoly} daje
$N(a)=(-1)^p(-a^p)=a^p$, ponieważ
$(-1)^{p+1}=1$ w ciele charakterystyki $p$
(dla $p=2$ znaki $1$ i $-1$ są identyczne).
W szczególności $N(s)=s^p=t\ne0$, choć
$\operatorname{Tr}_{L/K}(s)=0$.
\end{uwaga}

\section{Element pierwotny rozszerzenia rozdzielczego}
Przejście od wyznacznika do iloczynu wartości zanurzeń
wymaga rozdzielczości. Wyprowadzimy najpierw twierdzenie,
które pozwala opisać skończone rozszerzenie rozdzielcze
przez jeden generator.

\begin{twierdzenie}[Twierdzenie o elemencie pierwotnym]
\label{aa:tw:primitive-element}
Jeżeli $L/K$ jest skończonym rozszerzeniem rozdzielczym,
to istnieje $\gamma\in L$ takie, że $L=K(\gamma)$.
\end{twierdzenie}
\begin{proof}
Najpierw rozważmy przypadek nieskończonego ciała $K$.
Wystarczy udowodnić, że dla dowolnych $\alpha,\beta\in L$
istnieje $c\in K$ z $K(\alpha,\beta)=K(\alpha+c\beta)$;
powtarzanie tego kroku dla skończonej liczby generatorów
daje tezę.
Niech $\Omega$ będzie domknięciem algebraicznym $K$
zawierającym $L$; możemy je wybrać na mocy
Twierdzenia~\ref{aa:tw:algebraic-closure-exists}
i przedłużania zanurzeń. Niech
$\sigma_1,\ldots,\sigma_n:L\hookrightarrow\Omega$
będą wszystkimi $K$-zanurzeniami,
gdzie $n=[L:K]$ na mocy
Twierdzenia~\ref{aa:tw:embeddings-count}. Jeżeli
$(\sigma_i(\alpha),\sigma_i(\beta))\ne
(\sigma_j(\alpha),\sigma_j(\beta))$, to równość
$\sigma_i(\alpha+c\beta)=\sigma_j(\alpha+c\beta)$
może zachodzić dla co najwyżej jednej wartości $c\in K$:
różnica wartości jest liniowa w $c$, a jej współczynniki
nie są jednocześnie zerowe. Ponieważ par indeksów jest
skończenie wiele, wybieramy $c\in K$ omijające wszystkie
wartości powodujące równość dla par zanurzeń różniących się
na $K(\alpha,\beta)$. Wtedy liczba różnych obrazów
$\gamma=\alpha+c\beta$ przez $K$-zanurzenia $L$ jest równa
liczbie różnych ograniczeń tych zanurzeń do
$K(\alpha,\beta)$. Z twierdzenia o liczbie zanurzeń
i przedłużaniu zanurzeń ta liczba wynosi
$[K(\alpha,\beta):K]$. Tymczasem liczba różnych obrazów
$\gamma$ to $[K(\gamma):K]$, ponieważ $\gamma$ jest
rozdzielczy nad $K$. Mamy zatem
$[K(\gamma):K]=[K(\alpha,\beta):K]$.
Ponieważ $K(\gamma)\subseteq K(\alpha,\beta)$,
twierdzenie o wieży daje równość tych ciał.

Równość liczby obrazów z wymiarem $[K(\gamma):K]$
wynika z rozdzielczości: minimalny wielomian $\gamma$
ma dokładnie tyle różnych pierwiastków, ile wynosi jego stopień,
a każdy obraz $\gamma$ pod $K$-zanurzeniem jest jednym z nich.
Z drugiej strony każde $K$-zanurzenie $K(\gamma)$
przedłuża się do $K$-zanurzenia $L$, zatem żaden z tych
pierwiastków nie został pominięty. Analogiczne uzasadnienie
dotyczy ograniczeń do $K(\alpha,\beta)$.
Ponieważ $L/K$ jest skończone, wybór skończonego zbioru
generatorów $L$ nad $K$ nie wymaga dodatkowego twierdzenia:
wystarczy wziąć dowolną bazę $L$ jako przestrzeni wektorowej
i dołączać jej elementy po kolei.

Jeżeli $K$ jest skończone, to również $L$ jest ciałem
skończonym. Na mocy
Twierdzenia~\ref{aa:tw:finite-field-multiplicative-cyclic}
grupa $L^\times$ jest cykliczna. Niech $\gamma$ będzie
jej generatorem. Każdy niezerowy element $L$ jest potęgą
$\gamma$, więc $L=K(\gamma)$.
\end{proof}

\section{Ślad i norma jako sumy oraz iloczyny sprzężeń}
\begin{twierdzenie}[Wzory przez zanurzenia]
\label{aa:tw:trace-norm-embeddings}
Niech $L/K$ będzie skończonym rozszerzeniem rozdzielczym,
$n=[L:K]$, a $\Omega$ domknięciem algebraicznym $K$
zawierającym $L$. Niech $\sigma_1,\ldots,\sigma_n$
będą wszystkimi $K$-zanurzeniami $L\to\Omega$.
Wtedy dla każdego $a\in L$ zachodzą
\[
\operatorname{Tr}_{L/K}(a)=\sum_{i=1}^n\sigma_i(a),\qquad
N_{L/K}(a)=\prod_{i=1}^n\sigma_i(a).
\]
\end{twierdzenie}
\begin{proof}
Z Twierdzenia~\ref{aa:tw:primitive-element} wybieramy
$\gamma\in L$ z $L=K(\gamma)$. Niech $\gamma_i=\sigma_i(\gamma)$;
są one parami różne, ponieważ minimalny wielomian $\gamma$
jest rozdzielczy. W bazie $1,\gamma,\ldots,\gamma^{n-1}$
endomorfizm $m_a$ ma macierz $A=(a_{ij})$ o współczynnikach
w $K$. Dla każdego $j$ mamy
$a\gamma^{j-1}=\sum_i a_{ij}\gamma^{i-1}$.
Po zastosowaniu $\sigma_k$ dostajemy
$\sigma_k(a)\gamma_k^{j-1}=\sum_i a_{ij}\gamma_k^{i-1}$.
Jeśli $V=(\gamma_k^{j-1})_{k,j}$ jest macierzą Vandermonde'a,
to te równości dają
$VA=\operatorname{diag}(\sigma_1(a),\ldots,\sigma_n(a))V$.
Wyznacznik $V$ wynosi
$\prod_{i<j}(\gamma_j-\gamma_i)\ne0$, więc macierze
$A$ i diagonalna są podobne nad $\Omega$. Ich ślady
i wyznaczniki są równe, co daje obie żądane formuły.
\end{proof}

\begin{lemat}[Wyznacznik przekształcenia po ograniczeniu skalarów]
\label{aa:lem:restriction-determinant}
Niech $F/K$ będzie skończonym rozszerzeniem ciał stopnia $d$
i niech $r\ge1$ będzie liczbą całkowitą.
Niech $T:F^r\to F^r$ będzie dowolnym $F$-liniowym endomorfizmem.
Jeżeli $A\in M_r(F)$ jest jego macierzą w bazie standardowej,
to po potraktowaniu $F^r$ jako $K$-przestrzeni zachodzi
$\det_K T=N_{F/K}(\det_F A)$.
\end{lemat}
\begin{proof}
Jeżeli $A$ jest osobliwa, to $T$ ma niezerowy wektor
w jądrze również jako odwzorowanie $K$-liniowe.
Obie strony są wtedy zerowe.
Załóżmy, że $A$ jest odwracalna. Metoda eliminacji Gaussa
rozkłada $A$ na iloczyn macierzy elementarnych trzech rodzajów.
Wystarczy więc wykazać wzór dla każdej z nich,
ponieważ wyznaczniki oraz norma są multiplikatywne.

Jeśli macierz elementarna dodaje do jednej współrzędnej
wielokrotność innej, to po ograniczeniu skalarów
jej macierz jest blokową macierzą elementarną
z identycznościami na przekątnej i jednym blokiem
pozadiagonalnym. Jej wyznacznik nad $K$ wynosi $1$,
tak samo jak norma wyznacznika $1$ nad $F$.
Jeżeli mnożymy jedną współrzędną przez $u\in F^\times$,
to jako $K$-liniowy endomorfizm mamy na tej współrzędnej
mnożenie $m_u:F\to F$, a na pozostałych identyczność.
Wyznacznik nad $K$ jest zatem $N_{F/K}(u)$,
czyli normą wyznacznika nad $F$.
Wreszcie zamiana dwóch współrzędnych $F$ odpowiada
zamianie dwóch bloków o wymiarze $d$ nad $K$.
Jej wyznacznik wynosi $(-1)^{d^2}=(-1)^d$.
Wyznacznik tej operacji nad $F$ jest równy $-1$,
a $N_{F/K}(-1)=(-1)^d$ na mocy
Twierdzenia~\ref{aa:tw:trace-norm-basic}.
Wszystkie trzy rodzaje operacji spełniają wzór,
co kończy dowód.
\end{proof}

\begin{twierdzenie}[Przechodniość śladu i normy]
\label{aa:tw:trace-norm-tower}
\label{aa:thm:trace-norm-tower}
Niech $K\subseteq F\subseteq L$ będą ciałami takimi, że
$[F:K]<\infty$ i $[L:F]<\infty$
(nie wymagamy rozdzielczości).
Wtedy dla każdego $a\in L$ mamy
\[
\operatorname{Tr}_{L/K}(a)=\operatorname{Tr}_{F/K}\bigl(
\operatorname{Tr}_{L/F}(a)\bigr),\qquad
N_{L/K}(a)=N_{F/K}\bigl(N_{L/F}(a)\bigr).
\]
\end{twierdzenie}
\begin{proof}
Niech $r=[L:F]$, $d=[F:K]$ i wybierzmy bazę
$u_1,\ldots,u_r$ przestrzeni $L$ nad $F$.
Macierz $F$-liniowego mnożenia $m_a:L\to L$
oznaczmy przez $A=(a_{ij})\in M_r(F)$.
Jeśli $v_1,\ldots,v_d$ jest bazą $F$ nad $K$,
to $u_iv_j$ dla $1\le i\le r$, $1\le j\le d$
tworzą bazę $L$ nad $K$. W tej bazie macierz
$K$-liniowego mnożenia $m_a$ ma bloki
odpowiadające $K$-liniowym przekształceniom
$m_{a_{ij}}:F\to F$.
Zatem jej ślad to suma śladów bloków diagonalnych:
\[
\operatorname{Tr}_{L/K}(a)
=\sum_{i=1}^r\operatorname{Tr}_{F/K}(a_{ii})
=\operatorname{Tr}_{F/K}\bigl(\sum_{i=1}^r a_{ii}\bigr)
=\operatorname{Tr}_{F/K}(\operatorname{Tr}_{L/F}(a)).
\]
Dla normy stosujemy
Lemat~\ref{aa:lem:restriction-determinant}
do tego samego $F$-liniowego endomorfizmu $m_a$:
\[
N_{L/K}(a)=\det_K m_a
=N_{F/K}(\det_F A)=N_{F/K}(N_{L/F}(a)).
\]
Obie tożsamości są więc prawdziwe również dla
rozszerzeń nierozdzielczych.
\end{proof}

\begin{uwaga}[Wzory przez zanurzenia a rozdzielczość]
\label{aa:uw:trace-norm-inseparable}
Przechodniość śladu i normy jest twierdzeniem
czysto liniowoalgebraicznym, niewymagającym rozdzielczości.
Rozdzielczość jest natomiast konieczna w sformułowaniu
Twierdzenia~\ref{aa:tw:trace-norm-embeddings}
z dokładnie $[L:K]$ \emph{różnymi} zanurzeniami.
\end{uwaga}

\section{Forma śladu i wyróżnik}
\begin{definicja}[Forma śladu i wyróżnik bazy]
\label{aa:def:trace-discriminant}
Niech $L/K$ będzie skończonym rozszerzeniem stopnia $n$.
\emph{Formą śladu} na $L$ nazywamy $K$-dwuliniowe
odwzorowanie $B_{L/K}:L\times L\to K$,
$B_{L/K}(x,y)=\operatorname{Tr}_{L/K}(xy)$.
Dla bazy $b=(b_1,\ldots,b_n)$ definiujemy jej
\emph{wyróżnik}
$\operatorname{disc}_{L/K}(b)=\det\bigl(
\operatorname{Tr}_{L/K}(b_i b_j)\bigr)_{i,j=1}^n\in K$.
\end{definicja}

\begin{twierdzenie}[Niezdegenerowanie formy śladu]
\label{aa:tw:trace-nondegenerate}
Niech $L/K$ będzie skończonym rozszerzeniem rozdzielczym.
Wówczas forma śladu $B_{L/K}$ jest niezdegenerowana: dla każdego
$x\in L\setminus\{0\}$ istnieje $y\in L$ takie, że
$\operatorname{Tr}_{L/K}(xy)\ne0$. W szczególności
wyróżnik dowolnej bazy jest niezerowy.
\end{twierdzenie}
\begin{proof}
Niech $\Omega$ będzie domknięciem algebraicznym $K$
zawierającym $L$. Wybierzmy element pierwotny $\gamma$
i wszystkie $n=[L:K]$ różne zanurzenia
$\sigma_i:L\hookrightarrow\Omega$ na mocy
Twierdzenia~\ref{aa:tw:trace-norm-embeddings}. Dla bazy
$b_j=\gamma^{j-1}$ macierz formy śladu ma
współczynniki
$\operatorname{Tr}_{L/K}(b_jb_k)=\sum_i\sigma_i(b_j)\sigma_i(b_k)$.
Jeśli $V=(\sigma_i(b_j))$, to macierz formy wynosi
$V^{\mathsf T}V$. Zatem jej wyznacznik jest równy
$(\det V)^2$, a wyznacznik Vandermonde'a $V$ jest niezerowy.

Warto podkreślić, że macierz $V$ ma współczynniki w $\Omega$,
a niekoniecznie w $K$. Nie przeszkadza to w argumentacji:
niezerowość wyznacznika macierzy o współczynnikach w $K$
można sprawdzić po rozszerzeniu skalarów do $\Omega$.
W bazie potęgowej dostajemy ponadto konkretną formułę
\[
\operatorname{disc}_{L/K}(1,\gamma,\ldots,\gamma^{n-1})
=\prod_{1\le i<j\le n}(\sigma_j(\gamma)-\sigma_i(\gamma))^2.
\]
Prawa strona nie zależy od wybranej numeracji zanurzeń,
ponieważ zamiana dwóch wierszy macierzy $V$ zmienia
znak $\det V$, ale nie zmienia jego kwadratu.
Dla dowolnej innej bazy macierz formy zmienia się
na $S^{\mathsf T}(V^{\mathsf T}V)S$ dla odwracalnej
macierzy $S\in\operatorname{GL}_n(K)$, więc również
ma niezerowy wyznacznik. Macierz odwracalna oznacza,
że forma ma zerowe jądro, co daje tezę.
\end{proof}

\begin{przyklad}[Wyróżnik bazy kwadratowej]
\label{aa:prz:quadratic-discriminant}
Niech $d\in\mathbb Z\setminus\{0,1\}$ będzie liczbą
bezkwadratową, $L=\mathbb Q(\sqrt d)$, a
$b=(1,\sqrt d)$ bazą $L$ nad $\mathbb Q$.
Wówczas $\operatorname{Tr}(1)=2$,
$\operatorname{Tr}(\sqrt d)=0$ oraz
$\operatorname{Tr}(d)=2d$. Macierz formy śladu to
$\begin{psmallmatrix}2&0\\0&2d\end{psmallmatrix}$,
a wyróżnik jest równy $4d$. W szczególności
nie zanika, gdy $d\ne0$.
\end{przyklad}

\begin{lemat}[Zmiana bazy i wyróżnik]\label{aa:lem:trace-discriminant-base-change}
Niech $L/K$ będzie dowolnym skończonym rozszerzeniem
stopnia $n\ge1$ i niech $b=(b_1,\ldots,b_n)$,
$c=(c_1,\ldots,c_n)$ będą dwiema bazami $L$ nad $K$.
Jeżeli $S=(s_{ij})\in\operatorname{GL}_n(K)$
spełnia $c_j=\sum_{i=1}^n s_{ij}b_i$, to
\[
\operatorname{disc}_{L/K}(c)=(\det S)^2
\operatorname{disc}_{L/K}(b).
\]
Nie zakładamy rozdzielczości; obie strony mogą być zerowe.
\end{lemat}
\begin{proof}
Niech $G_b=(\operatorname{Tr}_{L/K}(b_ib_j))_{i,j}$.
Dla dowolnych $j,k$ dwuliniowość formy śladu daje
\[
\operatorname{Tr}_{L/K}(c_jc_k)
=\sum_{i=1}^n\sum_{h=1}^n
s_{ij}s_{hk}\operatorname{Tr}_{L/K}(b_ib_h).
\]
To właśnie wzór na wyraz $(j,k)$ macierzy
$G_c=S^{\mathsf T}G_bS$. Biorąc wyznaczniki,
otrzymujemy
\[
\det G_c=\det(S^{\mathsf T})\det G_b\det S
=(\det S)^2\det G_b,
\]
jak żądano.
\end{proof}
\begin{przyklad}[Wyróżnik zależy od bazy]
\label{aa:prz:trace-discriminant-other-basis}
Weźmy $L=\mathbb Q(\sqrt5)$. Dla bazy
$b=(1,\sqrt5)$ poprzedni przykład daje $\operatorname{disc}_{L/\mathbb Q}(b)=20$.
Rozważmy inną bazę $c=(1,(1+\sqrt5)/2)$.
Jej macierz przejścia z bazy $b$ wynosi
$S=\bigl(\begin{smallmatrix}1&1/2\\0&1/2\end{smallmatrix}\bigr)$,
więc $\det S=1/2\ne0$. Z lematu otrzymujemy
$\operatorname{disc}_{L/\mathbb Q}(c)=(1/2)^2\cdot20=5$.
Można to skontrolować bezpośrednio:
$\operatorname{Tr}(1)=2$,
$\operatorname{Tr}((1+\sqrt5)/2)=1$, a
$((1+\sqrt5)/2)^2=(3+\sqrt5)/2$, którego ślad wynosi $3$.
Macierz formy w bazie $c$ to
$\bigl(\begin{smallmatrix}2&1\\1&3\end{smallmatrix}\bigr)$
i jej wyznacznik wynosi $6-1=5$.
Zatem wyróżnik \emph{bazy} nie jest niezmiennikiem liczbowym
samego rozszerzenia; jego klasa modulo niezerowe kwadraty
jest niezależna od bazy. Nie należy tej definicji utożsamiać
bez dodatkowych założeń z wyróżnikiem wielomianu.
\end{przyklad}


\section{Ślad i norma w ciałach skończonych}
\begin{twierdzenie}[Jawne wzory nad ciałem pierwszym]
\label{aa:tw:finite-field-trace-norm}
Niech $q=p^n$, gdzie $p$ jest liczbą pierwszą,
a $n\ge1$. Dla $a\in\mathbb F_q$ zachodzą
\[
\operatorname{Tr}_{\mathbb F_q/\mathbb F_p}(a)
=\sum_{j=0}^{n-1}a^{p^j},\qquad
N_{\mathbb F_q/\mathbb F_p}(a)=a^{1+p+\cdots+p^{n-1}}.
\]

Wykładnik w formule normy wynosi
$1+p+\cdots+p^{n-1}=(p^n-1)/(p-1)$.
Dla $a=0$ powyższy wzór na normę również daje $0$,
ponieważ wykładnik jest dodatni (także dla $n=1$).
Oba odwzorowania są surjektywne na odpowiednio
$(\mathbb F_p,+)$ i $(\mathbb F_p^\times,\cdot)$,
przy czym dla normy ograniczamy dziedzinę do $\mathbb F_q^\times$.
\end{twierdzenie}
\begin{proof}
Grupę Galois $\mathbb F_q/\mathbb F_p$ generuje
$\operatorname{Fr}_p:a\mapsto a^p$ na mocy
Twierdzenia~\ref{aa:tw:finite-field-frobenius-group}.
Wzory wynikają z
Twierdzenia~\ref{aa:tw:trace-norm-embeddings}.
Ponieważ $\mathbb F_q^\times$ jest cykliczna rzędu $q-1$,
obraz homomorfizmu $x\mapsto x^{(q-1)/(p-1)}$
ma rząd
$(q-1)/\gcd(q-1,(q-1)/(p-1))=p-1$.
Jest to cała grupa $\mathbb F_p^\times$.
Ślad jest liniowym odwzorowaniem przestrzeni
$\mathbb F_q$ nad $\mathbb F_p$ w jednowymiarową
przestrzeń $\mathbb F_p$. Nie jest zerowy: w przeciwnym
razie każdy z $q$ elementów byłby pierwiastkiem
wielomianu $T(X)=\sum_{j=0}^{n-1}X^{p^j}$,
którego stopień wynosi $p^{n-1}<q$ i który jest
niezerowy. Wobec tego obraz śladu jest niezerową
podprzestrzenią jednowymiarowej przestrzeni $\mathbb F_p$,
a więc całym $\mathbb F_p$.
\end{proof}

\begin{przyklad}[Ślad w ciele czteroelementowym]
\label{aa:prz:finite-field-trace-four}
Niech $\mathbb F_4=\mathbb F_2(\omega)$, gdzie
$\omega^2+\omega+1=0$. Wówczas $\omega^2=\omega+1$
i $\omega^3=1$, więc $\mathbb F_4=\{0,1,\omega,\omega^2\}$.
Z poprzedniego twierdzenia dla $p=n=2$ mamy
$\operatorname{Tr}(x)=x+x^2$ i $N(x)=x^3$.
Dostajemy
\[
\begin{array}{c|cccc}
x&0&1&\omega&\omega^2\\\hline
\operatorname{Tr}(x)&0&0&1&1\\
N(x)&0&1&1&1
\end{array}
\]
Na przykład $\operatorname{Tr}(\omega)=\omega+\omega^2
=\omega+(\omega+1)=1$, a
$N(\omega)=\omega^{1+2}=\omega^3=1$.
Ten rachunek pokazuje, dlaczego w charakterystyce $p$
nie wystarczy sprawdzić $\operatorname{Tr}(1)$:
tutaj $\operatorname{Tr}(1)=2=0$, mimo że ślad jest surjektywny.
\end{przyklad}

\begin{uwaga}[Co dalej?]
\label{aa:uw:trace-norm-next}
Ślad i norma są użyteczne w arytmetyce ciał liczbowych,
teorii form kwadratowych i opisie rozszerzeń cyklicznych.
W następnym rozdziale połączymy grupy Galois z pytaniem,
czy pierwiastki wielomianu można wyrazić za pomocą
skończonego ciągu działań arytmetycznych i wyciągania pierwiastków.
\end{uwaga}
